\subsection{One-Time Signatures}
\label{sec:signatures}

Each input provider authenticates one ciphertext using an ordinary
one-time signature.  No deniability is needed: the ciphertext remains
fixed when the encryption simulator later explains the sender's randomness.

A signature scheme has polynomial-time algorithms
\[
  (\sigvk,\sigsk)\sample\SigGen(1^\lambda),\qquad
  \sig\sample\Sign(\sigsk,m),\qquad
  \Verify(\sigvk,m,\sig)\in\{\top,\bot\}.
\]
Messages have a prescribed polynomial length bound; shorter messages
are encoded and padded unambiguously.  Correctness requires an honestly
generated signature to verify, except with negligible probability.
Key-generation and signing coins are retained and disclosed if the
signer is corrupted.

\begin{definition}[One-time unforgeability]
\label{def:one-time-signatures}
Give an efficient adversary these commands:
\begin{description}
  \item[$\GetVerificationKey()$.]
    Generate $(\sigvk,\sigsk)$ and return $\sigvk$.
  \item[$\SignMessage(m)$.]
    Record $m$ and return $\Sign(\sigsk,m)$.
  \item[$\Forge(m^*,\sig^*)$.]
    Terminate.  The adversary wins if
    $\Verify(\sigvk,m^*,\sig^*)=\top$ and $m^*$ was not the
    message submitted to $\SignMessage$.
\end{description}
The first two commands are accepted at most once each, in that order.
The adversary may attempt its forgery after receiving the verification
key, whether or not it requested a signature.  Malformed commands are
rejected without changing state.  The scheme is one-time unforgeable
if every efficient adversary's winning probability is negligible.
\end{definition}

Lamport's construction~\cite{Lamport1979Signatures} suffices for bounded
messages from one-way functions.  Signing a collision-resistant digest
instead keeps the signature length independent of the ciphertext length;
this hash-then-sign variant additionally requires collision resistance.
We use a one-time signature scheme as a black box, and keep the signed
message's session, round, input-role, and sample identifiers explicit.

Unforgeability protects the message before signing-key exposure.  Once
the signed ciphertext is irrevocably recorded, the ledger fixes that
input even if the signing key later leaks.
